\paragraph{Kiểm toán thực tế và giới hạn nhãn}
Bộ dữ liệu thật có 180519 dòng, 53 cột và 65752 đơn hàng. Có 45902 đơn nhiều dòng, không có dòng trùng hoàn toàn và không có đơn có nhãn mâu thuẫn giữa các dòng. Số dòng lặp mã đơn là 114767; đây không phải bằng chứng dữ liệu trùng sai. Các đơn nhiều dòng được đặt trọng số lớn hơn trong metric theo dòng.
Nhãn gồm 98977 dòng dương và 81542 dòng âm. Đối chiếu Delivery Status cho thấy 7754 dòng ``Shipping canceled'' nằm ở lớp 0, cùng với 41592 dòng giao sớm và 32196 dòng đúng hạn. Vì vậy lớp 0 là ``không được gán giao trễ'' theo nguồn, không đồng nghĩa mọi đơn đã giao thành công đúng hạn. Chưa lọc đơn hủy để giữ đúng target gốc; việc đổi quần thể cần thỏa thuận và tái thiết kế thí nghiệm.
Missing chỉ xuất hiện ở Customer Lname (8), Customer Zipcode (3), Order Zipcode (155679) và Product Description (180519). Đây đều là trường bị loại; imputer vẫn cần cho dữ liệu triển khai. Không thể dùng kết quả này để kết luận median/mode cải thiện metric trên tập hiện tại, vì đặc trưng được chọn không có missing. Ngoại lệ IQR chỉ được đánh dấu, không bị xóa. Địa danh chứa chuỗi có khả năng lỗi mã hóa cần được rà soát; giải mã CP1252 thành công không đảm bảo mọi chuỗi đều đúng ngữ nghĩa.
\paragraph{Diễn giải mô hình và lỗi có bằng chứng}
Với cấu hình đã chọn tree B, độ quan trọng theo giảm impurity của Shipping Mode Standard Class là 0.5735, giờ tạo đơn là 0.1758 và Same Day là 0.1356. Đây là độ quan trọng nội tại, có thiên lệch theo cardinality và không phải tác động nhân quả. Quy tắc cây được xuất trong bảng minh chứng; ngưỡng số thuộc thang StandardScaler, không diễn giải trực tiếp thành giờ gốc.
Mô hình có 248 FN, đều thuộc Standard Class (248/254 dòng dương của nhóm, FNR 97.64\%). Nhóm này chỉ có 16 FP trên 450 dòng âm. First Class có 10 FP/10 dòng âm và Second Class 40 FP/40 dòng âm: cây gần như cảnh báo toàn bộ hai nhóm này. Quy tắc cây ưu tiên phương thức giao hàng giải thích được khuynh hướng dự đoán; chưa giải thích được nguyên nhân vận hành khiến đơn thật bị trễ.
Ví dụ FN có đơn vị dòng Standard Class, kế hoạch 4 ngày, thị trường Europe; ví dụ FP gồm Second Class, kế hoạch 2 ngày, LATAM. Các ví dụ cùng đặc điểm có thể thuộc một đơn: không coi chúng là bằng chứng độc lập. Báo cáo lỗi đã bỏ mã đơn/khách, tên, email, địa chỉ và thời điểm chính xác. Sau quan sát lỗi test, không tinh chỉnh lại trên test này.
Luồng B tăng CV F1 của cây từ 0.6741 lên 0.7058; test từ 0.6688 lên 0.7098. Với K-NN, CV tăng từ 0.6668 lên 0.6714; test từ 0.6768 lên 0.6938. Khác biệt chỉ là bằng chứng trên một split tập con; cần lặp đánh giá theo nhóm hoặc holdout thời gian độc lập. Dummy có Recall=1 và test F1=0.6991 do dự đoán mọi dòng là trễ, nhưng precision chỉ 0.5374; cây B tăng precision với đánh đổi recall còn 0.6076. Chưa chứng minh cây B tối ưu chi phí kinh doanh, vì chưa có chi phí FN/FP.
